% =====================================================================
\section{Evaluación y Conclusiones}
% =====================================================================

% ---------------------------------------------------------------------
\subsection{Resultados obtenidos}
% ---------------------------------------------------------------------

El período cerró con una línea base de requerimientos completa y trazada. Los 13 requerimientos de usuario y los 21 de sistema quedaron enlazados por 35 relaciones de trazabilidad, con cobertura completa en ambos sentidos: ningún requerimiento de usuario quedó sin traducción técnica y ninguno de sistema quedó sin justificación. Doce de ellos, clasificados con riesgo alto, cuentan con una mitigación definida, y la verificación futura quedó especificada en 19 ensayos.

La simulación confirmó que el enlace cierra. Con 1 W de potencia de transmisión y las antenas supuestas, la relación Eb/N0 supera el umbral de 14 dB en 4 a 5 pasadas diarias sobre territorio nacional, con ventanas útiles del orden de 10 minutos y una distancia media de enlace cercana a 1.500 km. Los máximos alcanzan valores en torno a 30 dB. La Figura \ref{img:downlinksim} muestra el comportamiento del enlace de bajada a lo largo de dos días de propagación orbital: cada pico corresponde a una pasada sobre la estación terrena y la línea segmentada indica el umbral.

La simulación reveló además un comportamiento que un cálculo con ganancia constante no habría detectado. En las pasadas de alta elevación, la relación Eb/N0 cae por debajo del umbral en el instante de máxima aproximación: 12,9 dB a 82,7$^{\circ}$ de elevación y 655 km de distancia, pese a que ese es el punto de menor pérdida de espacio libre. La causa es el nulo cenital del patrón de radiación del monopolo terrestre, visible en la Figura \ref{img:ventana}. La consecuencia operativa es concreta: la ventana de contacto no es continua, y la planificación de la descarga debe evitar el paso por el cénit.

Según lo planificado, el proyecto cerró con el cronograma que se presentó en la Figura \ref{img:gantt} y con una estimación de costos de orden de magnitud de 133.000 USD para entregar un modelo de vuelo integrado, desglosada en la Tabla \ref{tab:rom}.

\fig[downlinksim]{Presupuesto de enlace de bajada para dos días de propagación orbital. Elaboración propia.}{width=\textwidth}{crop_downlink.png}

\fig[ventana]{Detalle de una ventana de contacto individual. Elaboración propia.}{width=\textwidth}{crop_zoom_uplink.png}

\begin{table}[H]
\centering
\begin{threeparttable}
\caption{Estimación ROM para un modelo de vuelo integrado. Elaboración propia.\vspace{1ex}}
\label{tab:rom}
\small
\begin{tabularx}{\textwidth}{|@{\hspace{4pt}}l X r@{\hspace{4pt}}|}
\toprule
\textbf{Ítem} & \textbf{Descripción} & \textbf{Costo estimado (USD)} \\
\midrule
Recursos humanos          & Dos ingenieros durante 12 meses                & 48.000 \\
Componentes y materiales  & Electrónica, mecánica y materiales             & 50.000 \\
AIT                       & Integración, pruebas funcionales y ambientales & 15.000 \\
Logística                 & Transporte, envíos, importaciones y soporte    & 20.000 \\
\midrule
\textbf{Total}            &                                                & \textbf{133.000} \\
\bottomrule
\end{tabularx}
\begin{tablenotes}[flushleft]\footnotesize
\item Nota. La estimación asume el uso de componentes y materiales comerciales
disponibles en el mercado y la ejecución de las actividades de integración,
ensamblado y pruebas en instalaciones nacionales.
\end{tablenotes}
\end{threeparttable}
\end{table}


% ---------------------------------------------------------------------
\subsection{Análisis, limitaciones e implicancias futuras}
% ---------------------------------------------------------------------

Los resultados permiten afirmar que la arquitectura propuesta es viable a nivel preliminar. El enlace cierra con margen, y la restricción real de la misión no es la potencia sino el tiempo acumulado de contacto: entre 40 y 50 minutos diarios sobre territorio nacional. La arquitectura basada en radio definida por \textit{software} resuelve al mismo tiempo la tolerancia a la radiación, la modularidad y la actualización en vuelo, tres requerimientos que en una implementación con \textit{hardware} dedicado habrían competido entre sí.

El alcance de esas afirmaciones tiene límites explícitos. El modelo de enlace no incorpora pérdidas atmosféricas ni interferencia co-canal, y supone una sola estación remota transmitiendo, de modo que no representa las colisiones del escenario real. La preselección de componentes es preliminar y no verifica la dosis de radiación acumulada ni el presupuesto de potencia por modo operativo. Los 19 ensayos están especificados pero ninguno ejecutado: valores como la sensibilidad de $-115$ dBm son estimaciones de diseño. Y la estimación de costos, por definición, es de orden de magnitud.

Las decisiones tomadas proyectan además consecuencias sobre las fases siguientes. Adoptar un modelo MBSE obliga a mantenerlo actualizado. Elegir una plataforma SDR comercial crea dependencia de un proveedor único y de su cadena de suministro internacional. Los pasos a seguir están definidos: cerrar el presupuesto de potencia por modo y el análisis térmico, confirmar la selección de componentes, y ejecutar la campaña de ensayos en las instalaciones del Centro Espacial Nacional antes de la Revisión Crítica de Diseño.


% ---------------------------------------------------------------------
\subsection{Cumplimiento del objetivo de la empresa}
% ---------------------------------------------------------------------

Los dos objetivos que DTS definió para el período se cumplieron, aunque en distinto grado.

El primero, dar el primer paso formal en el desarrollo de capacidades espaciales, se cumplió y tuvo continuidad. La empresa cerró el período con un proyecto espacial formalizado y documentado bajo su propio sistema. Dos hechos posteriores lo confirman: DTS inició la capacitación del Departamento de I+D+i en desarrollo de proyectos espaciales, y BIFROST continúa en desarrollo tras el término de la práctica. El objetivo no quedó como un ejercicio aislado, sino que abrió efectivamente la línea de trabajo.

El segundo, aprovechar la ventana de transferencia tecnológica con INVAP, se cumplió de forma parcial. La contraparte participó en cuatro instancias formales de revisión y su retroalimentación quedó incorporada en los requerimientos y en la caracterización del canal. La transferencia en materia de integración, ensamblado y ensayos, en cambio, queda pendiente para las fases siguientes.

El entregable en que ambos objetivos convergían también se alcanzó: la Revisión de Diseño Preliminar se cerró el 27 de febrero de 2026, dentro del plazo de la práctica.


% ---------------------------------------------------------------------
\subsection{Impacto y relevancia de los resultados}
% ---------------------------------------------------------------------

El desarrollo de BIFROST permite responder la pregunta planteada al inicio de este informe. DTS pudo iniciar su línea espacial con un ingeniero de planta y un practicante porque separó la definición del marco de trabajo, aportada por la experiencia del supervisor, de su ejecución, apoyada en un estándar de revisiones y en una metodología de modelado que imponen estructura por construcción. Ninguna de las dos partes habría bastado por sí sola.

En lo técnico, la arquitectura propuesta es viable. El enlace cierra con margen sobre el umbral en 4 a 5 pasadas diarias, y la solución basada en radio definida por \textit{software} satisface los requerimientos críticos de tolerancia a la radiación, modularidad y actualización en vuelo.

El diseño, sin embargo, no está cerrado. Ningún ensayo se ha ejecutado, el escenario con múltiples estaciones remotas no está modelado y la selección de componentes espera confirmación. La Revisión Crítica de Diseño es el siguiente umbral.

Para la empresa, el resultado es que la línea espacial dejó de ser una declaración del plan estratégico y pasó a ser un proyecto en curso, con metodología propia, herramientas propias y un equipo en formación.
